\documentclass[12pt,a4paper,leqno]{article} % "leqno" pone la numeración a la izquierda
\usepackage{amsmath}
\usepackage[utf8]{inputenc}
\usepackage[T1]{fontenc}
\usepackage[spanish]{babel}

% Alinear ecuaciones totalmente a la izquierda
%\setlength{\mathindent}{0pt}

\begin{document}
	
	\begin{align}
		f(x) &= x^{4} - 4x^{2} \\
		f(x) &= x^{3} - 2x^{2} - 3x \\
		f(x) &= 6x^{4} + x^{3} - 15x^{2} \\
		f(x) &= x^{4} - 13x^{2} + 36 \\
		f(x) &= -x^{4} + 18x^{2} - 81 \\
		f(x) &= x^{3} - 3x^{2} - 9x + 27 \\
		f(x) &= 2x^{5} + 3x^{4} - 16x^{2} - 24x \\
		f(x) &= -x^{5} + 16x^{3} \\
		f(x) &= (x^{2} - 2x - 3)^{2} \\
		f(x) &= -\tfrac{1}{3}\left(2x^{4} + 9x^{3} - 5x^{2}\right)^{2}
	\end{align}
	
\end{document}


